% GENERATED FILE. Edit canonical lessons, then run npm run generate:books.
% canonical-source-hash: fnv1a64:0cb4601d3b1009dc
% canonical-lessons: GE-C525-streamen, GE-C525-abonnieren, GE-C525-videospiel, GE-C525-zocken, GE-C525-kopfhoerer

\chapter{Streamen und zocken}
\label{ch:ge-streamen-und-zocken}
\hlchaptermodality{525}

\begin{chapteropening}
\textbf{By the end of this chapter:} \emph{I can talk about streaming and video games.}

Complete the last lesson of chapter 525: I can talk about streaming and video games.
\end{chapteropening}

\section[streamen]{streamen — to stream}

\label{lesson:GE-C525-streamen}



\begin{quote}
Before the new one: say the German for subtitle, then the German for to dub.
\end{quote}

\subsection*{You'll want to know: streamen}
\textbf{streamen} — \textquotedblleft{}to stream\textquotedblright{}.

\begin{quote}
Wir streamen die meisten Filme.
\end{quote}

\begin{quote}
\emph{We stream most films.}
\end{quote}

Regular: \emph{hat gestreamt}. Said as in English.

\begin{cousinweb}
\emph{streamen} was borrowed from English \emph{to stream} (EtymDB).
\end{cousinweb}

\begin{grammarlens}[title={What you've built}]
One more word for telling, explaining and describing at B1.
\end{grammarlens}

\subsection*{Your turn}
\begin{itemize}
\raggedright
  \item \emph{Say it:} \emph{streamen}
  \item \emph{Say it:} \emph{streamen}, once more
  \item \emph{Say it:} a sentence of your own with \emph{streamen}
  \item \emph{Recall:} say the German for subtitle, then the German for to dub, then say \emph{streamen} again
\end{itemize}

\begin{tcolorbox}[breakable,colback=teal!4,colframe=teal!35!black,title={Before you move on}]
What does streamen mean? (\textquotedblleft{}To stream\textquotedblright{}.) Say it once more.
\end{tcolorbox}
\section[abonnieren]{abonnieren — to subscribe to}

\label{lesson:GE-C525-abonnieren}



\begin{quote}
Before the new one: say the German for to dub, then the German for to stream.
\end{quote}

\subsection*{You'll want to know: abonnieren}
\textbf{abonnieren} — \textquotedblleft{}to subscribe to\textquotedblright{}.

\begin{quote}
Ich habe eine Zeitschrift über Gärten abonniert.
\end{quote}

\begin{quote}
\emph{I've subscribed to a gardening magazine.}
\end{quote}

Regular, Perfekt without \emph{ge-}. The noun: \emph{das Abonnement}, in speech \emph{das Abo}.

\begin{cousinweb}
\emph{abonnieren} was borrowed from French \emph{abonner} \textquotedblleft{}to subscribe\textquotedblright{} (EtymDB).
\end{cousinweb}

\begin{grammarlens}[title={What you've built}]
One more word for telling, explaining and describing at B1.
\end{grammarlens}

\subsection*{Your turn}
\begin{itemize}
\raggedright
  \item \emph{Say it:} \emph{abonnieren}
  \item \emph{Say it:} \emph{abonnieren}, once more
  \item \emph{Say it:} a sentence of your own with \emph{abonnieren}
  \item \emph{Recall:} say the German for to dub, then the German for to stream, then say \emph{abonnieren} again
\end{itemize}

\begin{tcolorbox}[breakable,colback=teal!4,colframe=teal!35!black,title={Before you move on}]
What does abonnieren mean? (\textquotedblleft{}To subscribe to\textquotedblright{}.) Say it once more.
\end{tcolorbox}
\section[das Videospiel]{das Videospiel — video game}

\label{lesson:GE-C525-videospiel}



\begin{quote}
Before the new one: say the German for to stream, then the German for to subscribe to.
\end{quote}

\subsection*{You'll want to know: das Videospiel}
\textbf{das Videospiel} — \textquotedblleft{}video game\textquotedblright{}.

\begin{quote}
Zu Weihnachten wünscht sich mein Neffe ein neues Videospiel.
\end{quote}

\begin{quote}
\emph{My nephew wants a new video game for Christmas.}
\end{quote}

Plural: \emph{die Videospiele}.

\begin{cousinweb}
\emph{Videospiel} is built on \emph{Video}, borrowed from English \emph{video}, from Latin \emph{videō} \textquotedblleft{}I see\textquotedblright{}, and \emph{Spiel} \textquotedblleft{}game\textquotedblright{}, from Old High German \emph{spil} (EtymDB).
\end{cousinweb}

\begin{grammarlens}[title={What you've built}]
One more word for telling, explaining and describing at B1.
\end{grammarlens}

\subsection*{Your turn}
\begin{itemize}
\raggedright
  \item \emph{Say it:} \emph{das Videospiel}
  \item \emph{Say it:} \emph{das Videospiel}, once more
  \item \emph{Say it:} a sentence of your own with \emph{das Videospiel}
  \item \emph{Recall:} say the German for to stream, then the German for to subscribe to, then say \emph{das Videospiel} again
\end{itemize}

\begin{tcolorbox}[breakable,colback=teal!4,colframe=teal!35!black,title={Before you move on}]
What does das Videospiel mean? (\textquotedblleft{}Video game\textquotedblright{}.) Say it once more.
\end{tcolorbox}
\section[zocken]{zocken — to play video games; to gamble}

\label{lesson:GE-C525-zocken}



\begin{quote}
Before the new one: say the German for to subscribe to, then the German for video game.
\end{quote}

\subsection*{You'll want to know: zocken}
\textbf{zocken} — \textquotedblleft{}to play video games; to gamble\textquotedblright{}.

\begin{quote}
Die Jungen zocken den ganzen Nachmittag.
\end{quote}

\begin{quote}
\emph{The boys are gaming all afternoon.}
\end{quote}

Informal. It first meant to gamble.

\begin{cousinweb}
\emph{zocken} was borrowed from Yiddish \emph{tsokn} \textquotedblleft{}to gamble\textquotedblright{} (EtymDB), which is usually traced to Hebrew \emph{ṣāḥaq} \textquotedblleft{}to laugh, to play\textquotedblright{}.
\end{cousinweb}

\begin{grammarlens}[title={What you've built}]
One more word for telling, explaining and describing at B1.
\end{grammarlens}

\subsection*{Your turn}
\begin{itemize}
\raggedright
  \item \emph{Say it:} \emph{zocken}
  \item \emph{Say it:} \emph{zocken}, once more
  \item \emph{Say it:} a sentence of your own with \emph{zocken}
  \item \emph{Recall:} say the German for to subscribe to, then the German for video game, then say \emph{zocken} again
\end{itemize}

\begin{tcolorbox}[breakable,colback=teal!4,colframe=teal!35!black,title={Before you move on}]
What does zocken mean? (\textquotedblleft{}To play video games; to gamble\textquotedblright{}.) Say it once more.
\end{tcolorbox}
\section[der Kopfhörer]{der Kopfhörer — headphones}

\label{lesson:GE-C525-kopfhoerer}



\begin{quote}
Before the new one: say the German for video game, then the German for to play video games.
\end{quote}

\subsection*{You'll want to know: der Kopfhörer}
\textbf{der Kopfhörer} — \textquotedblleft{}headphones\textquotedblright{}.

\begin{quote}
Im Zug hört sie Musik mit Kopfhörern.
\end{quote}

\begin{quote}
\emph{On the train she listens to music with headphones.}
\end{quote}

Plural: \emph{die Kopfhörer}; one pair is often \emph{ein Kopfhörer}.

\begin{cousinweb}
\emph{Kopfhörer} is built on \emph{Kopf} \textquotedblleft{}head\textquotedblright{}, from Old High German \emph{kopf} \textquotedblleft{}cup\textquotedblright{}, Proto-Germanic \emph{kuppaz} (EtymDB), and \emph{Hörer} \textquotedblleft{}listener\textquotedblright{}, from \emph{hören} \textquotedblleft{}to hear\textquotedblright{}.
\end{cousinweb}

\begin{grammarlens}[title={What you've built}]
One more word for telling, explaining and describing at B1.
\end{grammarlens}

\subsection*{Your turn}
\begin{itemize}
\raggedright
  \item \emph{Say it:} \emph{der Kopfhörer}
  \item \emph{Say it:} \emph{der Kopfhörer}, once more
  \item \emph{Say it:} a sentence of your own with \emph{der Kopfhörer}
  \item \emph{Recall:} say the German for video game, then the German for to play video games, then say \emph{der Kopfhörer} again
\end{itemize}

\begin{tcolorbox}[breakable,colback=teal!4,colframe=teal!35!black,title={Before you move on}]
What does der Kopfhörer mean? (\textquotedblleft{}Headphones\textquotedblright{}.) Say it once more.
\end{tcolorbox}
